The relative displacement corresponding with plane strain is parallel to the
plane of the strain; or we have \(w=\) const., while \(u\) and \(v\) are
functions of \(x\) and \(y\) only. The axis of the resultant rotation is normal
to the plane of the strain. The cubical dilatation, \(\Delta\), and the
rotation, \(\varpi\), are connected with the displacement by the equations
\[
\Delta=\frac{\partial u}{\partial x}+\frac{\partial v^{\prime}}{\partial
y^{\prime}}, \quad 2 w=\frac{\partial v}{\partial x}-\frac{\partial u}{\partial
y}.
\]
We can have states of plane strain for which both \(\Delta\) and \(\varpi\)
vanish; the strain is pure shear, i.e. shearing strain combined with such a
rotation that the principal axes of the strain retain their primitive
directions. In any such state the displacement components \(v, u\) are
conjugate functions of \(x\) and \(y\), or \(v+\imath u\) is a function of the
complex variable \(x+\imath y\).

\section{Relations connecting the dilatation, the rotation and the displacement.} %15
The cubical dilatation \(\Delta\) is connected with the displacement
\((u, v, w)\) by the equation
\[
\Delta=\frac{\partial u}{\partial x}+\frac{\partial v}{\partial y}+
\frac{\partial w}{\partial z}.
\]
A scalar quantity derived from a vector by means of this formula is described
as the \textit{divergence} of the vector. We write
\begin{equation}\label{eq:14}
\Delta=\operatorname{div}(u, v, w).
\end{equation}
This relation is independent of coordinates, and may be expressed as
follows: - Let any closed surface \(S\) be drawn in the field of the vector, and
let \(N\) denote the projection of the vector on the normal drawn outwards at
any point on \(S\), also let \(d \tau\) denote any element of volume within
\(S\), then
\begin{equation}\label{eq:15}
\iint N d S=\iiint \Delta d \tau,
\end{equation}
the integration on the right-hand side being taken through the volume within
\(S\), and that on the left being taken over the
surface \(S\)\footnote{The result is a particular case of the theorem known as
'Green's theorem.' See Ency: \(d\). math. Wiss. II. A 2, Nos. 45-47.}.

The rotation \(\left(\varpi_x, \varpi_y, \varpi_z\right)\) is connected with
the displacement \((u, v, w)\) by the equations
\[
2 \varpi_x=\frac{\partial w}{\partial y}-\frac{\partial v}{\partial z}, \quad
2 \varpi_y=\frac{\partial u}{\partial z}-\frac{\partial w}{\partial x}, \quad
2 \varpi_z=\frac{\partial v}{\partial x}-\frac{\partial u}{\partial y} .
\]

A vector quantity derived from another vector by the process here indicated is
described as the curl of the other vector. We write
\begin{equation}\label{eq:16}
2\left(\varpi_x, \varpi_y, \varpi_z\right)=\operatorname{curl}(u, v, w).
\end{equation}
This relation is independent of coordinates\footnote{It is assumed that the
axes of \(x, y, z\) form a right-handed system.  If a transformation to a
left-handed system is admitted a convention must be made as to the sign of
the curl of a vector.}, and may be expressed as
follows:-Let any closed curve \(s\) be drawn in the field of the vector, and
let any surface \(S\) be described so as to have the curve \(s\) for an edge;
let \(T\) be the resolved part of the vector \((u, v, w)\) along the tangent at
any point of \(s\), and let \(2 \varpi_\nu\) be the projection of the vector
\(2\left(\varpi_x, \varpi_y, \varpi_z\right)\) on the normal at any point of
\(S\), then
\begin{equation}\label{eq:17}
\int T d s=\iint 2 \varpi_\nu d S,
\end{equation}
